\documentclass[11pt]{article}
\usepackage[margin=1in]{geometry}
\usepackage{amsmath,amssymb}
\usepackage{booktabs}
\usepackage{graphicx}
\usepackage{siunitx}
\usepackage[hidelinks]{hyperref}
\usepackage{xcolor}

\newcommand{\vdiff}{v_{\mathrm{diff}}}
\newcommand{\Dch}{\Delta}

\title{Identifiability of the aDDM, Encoder-Representation Comparison,\\
and Static Design Optimisation}
\author{CogFormer-2}
\date{2 September 2026}

\begin{document}
\maketitle

\begin{abstract}
Three connected results. First, the attentional drift-diffusion model (aDDM) has a
structural flat direction in $(v, \vdiff, \omega)$ when gaze is unobserved: the data
depend on these three parameters through only two functionals. Removing it by pinning
$v$ makes the inference target honest but does not recover $\omega$; letting item
values vary by trial recovers part of it (posterior contraction $0.013 \to 0.232$) but
runs into a second-order ceiling. Second, a paired comparison on shared validation data
overturns an earlier conclusion that a frozen transferred CogFormer summary is
statistically indistinguishable from a task-specific learned one: the two differ
clearly, the learned summary is better, and the gap roughly doubles as the design gets
richer. Third, because CogFormer amortises over designs, a static design-optimisation
sweep costs only forward passes; it shows that dichotomising a continuous predictor
carries an information penalty that \emph{grows} with sample size, equivalent to
needing $1.2$--$1.5\times$ the trials.
\end{abstract}

\section{The aDDM flat direction}

\subsection{Structure}

In the aDDM the relative accumulator's drift alternates with gaze between
$v_0 - \omega v_1$ (fixating item 0) and $\omega v_0 - v_1$ (fixating item 1), with
$v_0 = v$ and $v_1 = v + \vdiff$. In the fast-switching regime the data see this
telegraph process through two functionals only,
\begin{align}
  \mu     &= \tfrac{1}{2}(1+\omega)\,\vdiff, \label{eq:mu}\\
  \Dch    &= (1-\omega)\,(2v + \vdiff), \label{eq:delta}
\end{align}
where $\mu$ is the effective mean drift and $\Dch$ the attentional swing, which drives
both the first-fixation choice bias and the extra within-trial drift variance. Three
parameters mapping onto two functionals leaves exactly one flat direction, and $v$
spans it: $v$ never enters $\mu$ at all.

\subsection{Numerical confirmation}

Table~\ref{tab:flat} matches $\mu$ and $\Dch$ across triples with widely different
$v$ and $\omega$, simulating \num{200000} trials at $a=1$, $\tau=0.3$,
$s_v = s_\tau = 0$, \texttt{first\_fixation=item0}. The parameter $v$ moves by a
factor of $3.6$ and $\omega$ triples, yet choice proportion and RT quantiles are
indistinguishable. A control matching $\mu$ but not $\Dch$ shifts $P(c{=}1)$ by
$0.037$, confirming the flat direction is specifically the $(\mu,\Dch)$ level set and
not a general insensitivity.

\begin{table}[h]
\centering
\begin{tabular}{lll ccc}
\toprule
$v$ & $\vdiff$ & $\omega$ & $P(c{=}1)$ & RT $q_{10}/q_{50}/q_{90}$ \\
\midrule
1.00 & 0.500 & 0.25 & 0.5713 & 0.554 / 1.007 / 2.173 \\
1.67 & 0.417 & 0.50 & 0.5719 & 0.554 / 1.006 / 2.167 \\
3.57 & 0.357 & 0.75 & 0.5727 & 0.553 / 1.006 / 2.172 \\
\bottomrule
\end{tabular}
\caption{Three parameter triples matched on $(\mu,\Dch)$ produce
indistinguishable data.}
\label{tab:flat}
\end{table}

\section{Two attempted fixes}

Both arms were trained at an identical budget (500 epochs $\times$ 100 steps $\times$
batch 64, i.e.\ \num{3.2}M simulated datasets) and validated at 500 datasets
$\times$ 500 posterior draws. All contraction and correlation figures below are on the
\emph{natural} parameter scale, i.e.\ after applying the link functions; the saved
arrays are on the unconstrained scale, where a saturating link such as
$\omega$'s makes contraction misleading.

\subsection{Pinning $v$}

Fixing $v$ to its prior median ($1.0375 = \mathrm{softplus}(0.6)$) removes the flat
direction exactly: with $v$ known, $(\vdiff,\omega) \leftrightarrow (\mu,\Dch)$ is a
bijection. The inference target drops from 7 active cells to 6.

The reparameterisation succeeds structurally and fails practically
(Table~\ref{tab:vfixed}). $\omega$ reaches contraction $0.013$: its posterior is
$99.3\%$ as wide as its prior. Nothing else moves relative to the 7-dimensional run.

\begin{table}[h]
\centering
\begin{tabular}{l cc cc}
\toprule
& \multicolumn{2}{c}{frozen} & \multicolumn{2}{c}{baseline} \\
\cmidrule(lr){2-3}\cmidrule(lr){4-5}
cell & contraction & $r$ & contraction & $r$ \\
\midrule
$\vdiff$   & 0.939 & 0.971 & 0.939 & 0.970 \\
$\omega$   & \textbf{0.013} & 0.250 & \textbf{0.046} & 0.181 \\
$a$        & 0.983 & 0.992 & 0.984 & 0.991 \\
$\tau$     & 0.962 & 0.981 & 0.974 & 0.988 \\
$s_v$      & 0.101 & 0.332 & 0.033 & 0.219 \\
$s_\tau$   & 0.816 & 0.908 & 0.848 & 0.920 \\
\bottomrule
\end{tabular}
\caption{$v$-fixed target. Calibration holds throughout
($\mathrm{sd}(z)\in[0.95,1.02]$, 90\% coverage $\in[0.89,0.92]$), so this is a
likelihood property, not a training failure.}
\label{tab:vfixed}
\end{table}

The reason is that with $\vdiff$ free at inference time it absorbs any movement in
$\mu$. Holding $\mu$ \emph{constant} and varying $\omega$ from $0.1$ to $0.9$ moves
$P(c{=}1)$ by only $0.084$ with $s_v,s_\tau$ at their prior medians --- about
$3.7$ standard errors at 500 trials, reaching the data second-order through the
first-fixation transient's effect on RT shape, and readily absorbed by free
$a$, $\tau$, $s_v$, $s_\tau$. Removing the \emph{structural} flat direction left a
\emph{practical} one.

\subsection{Trial-varying item values}

The second design lets item values vary per trial and enter as observed continuous
regressors, $\vdiff^{(t)} = \beta_{d0} + \beta_{d1} u_1^{(t)}$ and
$v^{(t)} = \beta_{v0} + \beta_{v1} u_2^{(t)}$. Now $\vdiff$ enters both
Eq.~\eqref{eq:mu} and Eq.~\eqref{eq:delta} with the same $\beta_{d1}$, and the ratio of
its $\Dch$-effect to its $\mu$-effect is $2(1-\omega)/(1+\omega)$, free of
$\beta_{d1}$. So $\omega$ is structurally identified \emph{without} pinning $v$, and
$v$ returns to the target (9 active cells).

\begin{table}[h]
\centering
\begin{tabular}{l cc cc}
\toprule
& \multicolumn{2}{c}{frozen} & \multicolumn{2}{c}{baseline} \\
\cmidrule(lr){2-3}\cmidrule(lr){4-5}
cell & contraction & $r$ & contraction & $r$ \\
\midrule
$v$              & 0.126 & 0.188 & 0.154 & 0.140 \\
$\vdiff$         & 0.943 & 0.969 & 0.939 & 0.969 \\
$\omega$         & \textbf{0.232} & \textbf{0.488} & \textbf{0.319} & \textbf{0.609} \\
$a$              & 0.979 & 0.991 & 0.982 & 0.992 \\
$\tau$           & 0.975 & 0.987 & 0.976 & 0.989 \\
$s_v$            & 0.093 & 0.282 & 0.017 & 0.290 \\
$s_\tau$         & 0.893 & 0.943 & 0.868 & 0.939 \\
$\vdiff@r_1$     & 0.959 & 0.978 & 0.968 & 0.982 \\
$v@r_4$          & 0.322 & 0.553 & 0.453 & 0.654 \\
\bottomrule
\end{tabular}
\caption{Trial-varying values. $\omega$ improves roughly eighteen-fold in contraction
over the $v$-fixed target but remains weakly informed.}
\label{tab:varying}
\end{table}

\subsection{Where $\omega$'s information lives}

Table~\ref{tab:varying} localises the ceiling. Every quantity reaching the data through
$\mu$ is recovered cleanly ($\vdiff$ at $0.943$, its slope $\vdiff@r_1$ at $0.959$);
every quantity reaching it through $\Dch$ is not ($v$ at $0.126$, its slope $v@r_4$ at
$0.322$, $\omega$ at $0.232$). Since $\omega$ is identified only through $\Dch$, the
second-order nature of that channel is a genuine ceiling rather than an artefact of the
intercept-only design. Three independent lines of evidence agree on this: the direct
$\mu$-constant simulation, the $v$-fixed result, and the internal $\mu$/$\Dch$ split
here.

More trials will not rescue $\omega$: the signal grows as $\sqrt{n}$, so reaching
$r \approx 0.8$ would need roughly $10$--$16\times$ the data, i.e.\ 5000--8000 trials
per participant against the $\sim$300 typical of value-based choice experiments. The
remaining levers are observing fixations, which would put $\omega$ in the drift
first-order, or richer value manipulations.

\section{Comparing encoder representations}

\subsection{Why the earlier comparison was underpowered}

An earlier analysis concluded that a frozen CogFormer encoder (trained on DDM, RDM and
CDM, never on aDDM) and a task-specific BayesFlow \texttt{SetTransformer} were
statistically indistinguishable. That comparison used \emph{per-parameter marginal}
metrics computed on \emph{independently simulated} validation sets for each arm. Both
choices cost power, and the conclusion does not survive removing them.

\subsection{Method}

We simulate one validation set and push it through both trained approximators. Three
methodological points matter.

\paragraph{Pooled C2ST is degenerate here.} A classifier two-sample test that pools
posterior draws across datasets is uninformative for this comparison. Both arms are
calibrated, and a calibrated posterior averaged over the prior predictive returns the
\emph{prior}; the pooled sets are therefore both approximately the prior and the test
returns $0.5$ by construction, however much the per-dataset posteriors disagree. Only
the per-dataset statistic carries information.

\paragraph{The null is not $0.5$.} C2ST accuracy is biased upward in finite samples
with a flexible classifier. We construct an empirical null by running the identical
code path on two \emph{independent sampling calls from the same arm} --- same
dimension, same draw count, same classifier, same folds --- so the null carries the
identical finite-sample bias and only the excess over it is evidence. Measured nulls
ran $0.499$--$0.52$ depending on draw count. Arms and null share datasets, so they are
compared \emph{paired}.

\paragraph{Two proper scoring rules, for direction.} C2ST is symmetric and cannot rank
the arms. We use the expected log posterior density at the true parameter, which the
approximator provides exactly (the adapter's log-determinant Jacobian is folded in, so
densities are in the original parameter space), and CRPS, which is computed from draws
alone and is therefore immune to any constant offset between the arms' standardisers.

\subsection{Results}

\begin{table}[h]
\centering
\begin{tabular}{l cc}
\toprule
statistic & $v$-fixed (6 cells) & varying values (9 cells) \\
\midrule
C2ST joint, excess over null      & $+0.1176$ ($z{=}{+}29.7$) & $\mathbf{+0.2051}$ ($z{=}{+}38.0$) \\
C2ST identified cells only        & $+0.1200$ ($z{=}{+}30.3$) & $+0.1720$ ($z{=}{+}33.9$) \\
$\log q(\theta_{\text{true}}\mid x)$, frozen $-$ baseline & $-0.304$ ($z{=}{-}7.5$) & $\mathbf{-0.580}$ ($z{=}{-}10.4$) \\
CRPS total, frozen $-$ baseline   & $+0.0257$ ($+2.0\%$, $z{=}{+}4.5$) & $\mathbf{+0.1113}$ ($+5.3\%$, $z{=}{+}7.6$) \\
\bottomrule
\end{tabular}
\caption{Shared-validation comparison, 500 datasets $\times$ 500 draws. Positive CRPS
difference and negative log-density difference both favour the baseline. C2ST exceeds
its null on 100\% of datasets in both designs.}
\label{tab:arms}
\end{table}

The arms differ unambiguously, the learned task-specific summary is better on every
statistic, and the gap roughly doubles in relative terms as the design gets richer.

The per-cell CRPS breakdown on the trial-varying design separates two distinct
mechanisms. The learned summary wins where signal is \emph{marginal} --- $v@r_4$
($z{=}{+}4.9$), $\omega$ ($z{=}{+}3.6$) --- and, less expectedly, on the regression
slope $\vdiff@r_1$ ($z{=}{+}6.8$) despite that cell being well identified in both arms
($0.956$ vs $0.966$). The frozen summary is therefore worse not only where the signal
is weak but where the design carries \emph{regression structure the encoder never
trained on}: it saw DDM, RDM and CDM designs, never aDDM's value regressors.

\subsection{Revised claim}

The defensible statement is narrower than the original. On strongly identified
parameters the two summaries are indistinguishable ($\vdiff$, $a$, $\tau$ agree to
three decimals). The learned summary wins on marginal signal and on unfamiliar
regression structure, by roughly $2\%$ CRPS on a simple design and $5\%$ on a richer
one.

\section{Static design optimisation}

\subsection{Setup}

Because the design matrix is a \emph{conditioning input} rather than something baked
into training, a design sweep costs only forward passes on an existing checkpoint. A
per-design NPE or per-design MCMC would require one fit per cell.

The scenario: a psychologist has two predictors and must decide, for each, between
measuring it continuously (C) and reducing it to two levels (B). Both predictors load
on drift rate and threshold, giving 11 active cells (7 intercepts, plus
$(v,a)\times(u_1,u_2)$). Four designs --- CC, CB, BC, BB --- cross five trial counts.
The sweep stays inside $[200,1000]$, the range the checkpoint was trained on; below
that the curve would measure encoder degradation out of distribution rather than
information.

\subsection{The measure, and two traps}

The criterion is expected information gain, estimated as a Gaussian log-volume ratio,
\begin{equation}
  \mathrm{EIG} = \tfrac{1}{2}\log\det\Sigma_{\text{prior}}
               - \tfrac{1}{2}\,\mathbb{E}_x\!\left[\log\det\Sigma_{\text{post}}(x)\right].
\end{equation}

\paragraph{Matching dimension is necessary but not sufficient.} A $K$-level categorical
regressor contributes $K-1$ design columns, so for $K>2$ the parameter vector changes
dimension between designs and EIG becomes a log-volume in a different-dimensional
space. Restricting to $K=2$ fixes that --- the parameter mask is then identical across
all four designs, which the code asserts. But two further mismatches remain and were
measured directly: under the \texttt{cell\_means} coefficient prior a categorical
contrast has prior variance $2$ against $1$ for a continuous slope, and continuous
design columns are standardised to unit SD whereas a binary column has
$\mathrm{sd}=\sqrt{p(1-p)}$, measured at $0.33$ and $0.49$. Slope cells are therefore
rescaled by the realised SD of their own design column, putting every design on the
common estimand ``change in the linear predictor per SD of the predictor as run''.

\paragraph{Only the ratio is scale-free.} EIG, being a log-volume ratio, cancels that
rescaling and any fixed linear reparameterisation. The absolute posterior log-volume
does not, and ranks the all-binary design most certain purely because rescaling by
$0.33$ shrinks its estimand. EIG is the primary metric for this reason.

\subsection{Results}

\begin{table}[h]
\centering
\begin{tabular}{r cccc c}
\toprule
$n$ & CC & CB & BC & BB & CC $-$ BB \\
\midrule
200  & 17.418 & 16.568 & 16.562 & 16.308 & 1.110 \\
400  & 21.234 & 20.117 & 20.253 & 19.272 & 1.963 \\
600  & 23.540 & 21.748 & 22.145 & 21.426 & 2.115 \\
800  & 24.781 & 22.254 & 23.206 & 22.228 & 2.554 \\
1000 & 25.620 & 23.326 & 23.607 & 22.893 & 2.727 \\
\bottomrule
\end{tabular}
\caption{Expected information gain in nats. Standard errors are $0.11$--$0.18$.}
\label{tab:eig}
\end{table}

\begin{figure}[h]
\centering
\includegraphics[width=\textwidth]{../outputs/data/design_opt/ddm/design_eig.pdf}
\caption{Information gain and mean posterior contraction against trial count for the
four designs.}
\label{fig:eig}
\end{figure}

The ordering is $\mathrm{CC} > \mathrm{BC} \approx \mathrm{CB} > \mathrm{BB}$ at every
trial count, and the penalty for dichotomising \emph{grows} with sample size rather
than washing out.

\paragraph{Separating the trend from the prior.} The designs carry different priors
over the common estimand, worth $1.704$ nats between CC and BB. This is a constant in
$n$: it shifts the level of every gap but cancels exactly from the slope. Measuring
each design's EIG \emph{gain over its own $n{=}200$ baseline} removes the prior term,
and gives the trial count BB needs to match CC:

\begin{center}
\begin{tabular}{lcc}
\toprule
CC at & BB needs & multiplier \\
\midrule
$n=400$  & $\approx 479$  & $1.20\times$ \\
$n=600$  & $\approx 861$  & $1.43\times$ \\
$n=800$  & $\approx 1234$ & $1.54\times$ \\
$n=1000$ & $\approx 1486$ & $1.49\times$ \\
\bottomrule
\end{tabular}
\end{center}

Dichotomising both predictors costs $20$--$55\%$ more trials for the same information,
the multiplier growing before settling near $1.5\times$. This figure is prior-free.

\paragraph{Mechanism.} Per-200-trial EIG increments run $3.82, 2.31, 1.24, 0.84$ for CC
and $2.96, 2.16, 0.80, 0.67$ for BB. Both decelerate, but BB saturates faster: a binary
predictor has two design points, so additional trials re-sample the same cells, whereas
a continuous predictor keeps adding new ones. That is why the gap widens instead of
closing.

\section{Limitations}

\begin{enumerate}
\item \textbf{Absolute EIG gaps are partly prior.} Of the $1.11$--$2.73$ nat CC--BB
gap, $1.70$ is the \texttt{cell\_means} prior difference. At $n{=}200$ the raw gap is
smaller than that offset, so the data-driven component there favours BB; by $n{=}1000$
it favours CC by roughly $1.0$ nats. Only the trend and the trial multipliers are
clean. Per-design $\log\det\Sigma_{\text{prior}}$ is reported so this is checkable.

\item \textbf{This is factor-versus-parametric, not median-split.} Category proportions
are Dirichlet-drawn; realised split balance ran $0.38$--$0.45$ rather than $0.5$. The
literal median-split question needs balanced categories and arguably an ordinal design
generator, which is off-distribution for this checkpoint.

\item \textbf{$s_v$ is not recovered anywhere} and is not expected to be. It is the
classical second-order diffusion parameter, not an aDDM pathology.

\item \textbf{$\omega$'s identification is welded to the gaze-onset rule.} All results
use a deterministic first fixation. Under the standard random rule the first-fixation
advantage averages out and $\omega$'s already marginal signal would weaken further.

\item \textbf{The arm comparison is at matched training budget, not matched capacity.}
Each arm runs at its library's recommended defaults on the same number of simulated
datasets.
\end{enumerate}

\end{document}
